\documentclass[11pt]{article}
\usepackage[margin=1in]{geometry}
\usepackage{amsmath,amssymb,amsthm}
\usepackage{enumitem}
\usepackage[hidelinks]{hyperref}
\setlist{itemsep=2pt,topsep=3pt}
\newtheorem*{problem}{Open Problem (CP)}
\title{Erd\H{o}s Problem \#644: the remaining obstacle\\[2pt]\large A brief for reviewers}
\author{erdos644 project}
\date{November 2026}
\begin{document}
\maketitle

\section{The question}
Let $f(k,7)$ be the largest transversal number $\tau(H)$ of a $k$-uniform family $H$ in which every at most seven
members can be pierced by two points. Erd\H{o}s asked whether $f(k,7)=(\tfrac34+o(1))k$. The best published bounds are
$f(k,7)\le\lceil 7k/8\rceil$ and $f(4m,7)\ge 3m+1$ for $m\ge 10$ (Fon-Der-Flaass, Kostochka and Woodall, 1999).
We have reduced the question to one statement about covering families, plus two small residues. We ask for help with
that statement.

\section{The reduction (independently verified)}
\textbf{Game form.} If $\tau(H)>T$, every $T$-set misses some member. So a requester who names $T$-sets always receives
members avoiding them. If she asks every $T$-set, the answers on that single play form a family with $\tau>T$. Thus
$f(k,7)\le T$ if and only if she can always collect seven answers with no 2-point transversal.

\textbf{Confinement.} Three answers with empty common intersection confine all later analysis to their union $V$,
with $|V|\le 3k$.

\textbf{Co-edge form.} We work in units of $k$: $t=T/k=\tfrac34+\varepsilon$ and $N=|V|/k$. For a member $B$ its
\emph{co-edge} is $K=V\setminus B$, of mass $N-1$. Then:
\begin{itemize}
\item $\tau>T$ if and only if the co-edge family is \emph{$T$-covering}: every $T$-subset of $V$ lies in a co-edge.
\item Seven members have no 2-point transversal if and only if their co-edges cover every pair of $V$. We call such
seven co-edges a \emph{good seven}.
\end{itemize}
It suffices to treat all sufficiently small $\varepsilon>0$. We write $\omega=N-1-t$ and $\delta=N-\tfrac74$.

\textbf{Case split by $N$.} Every region is closed except one:
\begin{itemize}
\item $N\le 3t-\tfrac12$ (crowded): closed.
\item $3t-\tfrac12<N\le 5t-2$ (bottom strip): closed by a Frankl--R\"odl density argument.
\item $N\ge 3-t$ (petal side): closed except two small residues.
\item $5t-2<N<3-t$: the \textbf{core window}, open.
\end{itemize}

\section{The open problem}
\begin{problem}
Let $\varepsilon>0$ be sufficiently small, $t=\tfrac34+\varepsilon$, and $5t-2<N<3-t$ (about $1.75k$ to $2.25k$
points). Show that every $T$-covering family of co-edges contains a good seven.
\end{problem}

\noindent\textbf{A sufficient condition} (verified). The family contains one of:
\begin{enumerate}[label=(\alph*)]
\item a Q$'$ quadruple: four members with $\sigma\le(3t-N)/2$, where $\sigma$ is the mass lying in at most one of
them, and, for each of the three matchings of the four, (mass lying in exactly the pairs of that matching)
$+\,\sigma\le t$;
\item a pencil triple: three members whose union is $V$ and whose triple intersection is $\ge N-2t$;
\item a Theorem-T triangle: three members with empty triple intersection and pairwise intersections $a,b,d\le
(t+N-2)/2$ with $a+b+d\ge(3N-3-t)/2$.
\end{enumerate}
(Exact computation shows these three are not the only way to finish: about a third of forced good sevens contain none
of them. Every window family examined, however, contains at least one.)

\section{What a counterexample would have to look like (verified)}
A family with no good seven in the window cannot be any of the following:
\begin{itemize}
\item a bounded blow-up: membership depending only on profiles of at most $4\varepsilon k-12$ cells;
\item a fixed-order rule;
\item a rule choosing among priority orders by a window of at most $T-k/2$ points;
\item a family with kernel density above $2/3$;
\item a random or 6-wise independent selector with fewer than about $e^{1.19k}$ classes.
\end{itemize}
It must therefore be \emph{deterministic and strongly anti-correlated}. A random $T$-covering family of the relevant
size contains $e^{\Theta(k)}$ good sevens. At the smallest genuine window instance $(k,T,n)=(10,8,21)$, every
cell-symmetric search is UNSAT with checked certificates. The most anti-correlated family we know is pencil-free at
$\varepsilon=0$, but it fails there.

\section{What does not work}
Twenty-one approaches were analysed, each with an identified obstruction (details in \texttt{FAILED\_ROUTES.pdf}):
\begin{itemize}
\item bounded-request games;
\item junta approximation, because the families have density $e^{-\gamma k}$;
\item spread and sunflower lemmas, because the families are at most $2.3$-spread;
\item fractional relaxations of both the transversal and the pair-cover LP;
\item hypergraph containers;
\item flag algebras;
\item forbidden-intersection density arguments above the bottom strip;
\item adaptive restriction to links, since links lower $\varepsilon$ while keeping $\omega$;
\item entropy methods, since certificates need cyclic constraints;
\item covering-design stability;
\item several reformulations.
\end{itemize}
The lesson is that a successful argument must use the \emph{full} covering hypothesis, not only density, together with
\emph{adaptive} choice. It must also handle near-pool families, which sit exactly at the static slack $4\varepsilon$.

\section{What we would value}
\begin{enumerate}
\item Any tool that combines a covering hypothesis (``every $T$-set lies in a member'') with avoidance of a fixed small
configuration, in a linear, exponentially sparse regime.
\item A reason to believe the answer is \emph{no}: an idea for a deterministic, anti-correlated $T$-covering family with
no good seven.
\item Errors in any claim marked verified.
\item References to related exact results or techniques.
\end{enumerate}

\medskip\noindent Full notes, code and verification reports:
\url{https://github.com/cartiyoungercousin/erdos644}. Start with \texttt{README.md}; \texttt{docs/RESULTS.md} lists
every result with its verification.
\end{document}
